
\documentclass[aspectratio=169,10pt]{beamer}

\usetheme{Madrid}
\usecolortheme{default}

\setbeamercolor{palette primary}{bg=blue!80!black, fg=white}
\setbeamercolor{palette secondary}{bg=blue!60!black, fg=white}
\setbeamercolor{palette tertiary}{bg=blue!40!black, fg=white}
\setbeamercolor{palette quaternary}{bg=blue!20!black, fg=white}

\setbeamercolor{structure}{fg=blue!70!black}
\setbeamercolor{block title}{bg=blue!75!black, fg=white}
\setbeamercolor{block body}{bg=blue!5!white}
\setbeamercolor{block title alerted}{bg=red!75!black, fg=white}

% -----------------------------------------------------------
\begin{document}

% -----------------------------------------------------------
\section{State of the Art}
% ============================================================

% ------------------------------------------------------------
\begin{frame}{Paper MERLIN: Simulasi Spektrum Fisika Plasma}

\begin{columns}[T]

\begin{column}{0.56\textwidth}

\begin{block}{Identitas Artikel}
\small
	extbf{Judul}: \textit{MERLIN, an adaptive LTE radiative transfer model for any mixture: Validation on Eurofer97 in argon atmosphere}\\[4pt]

	extbf{Jurnal}: Journal of Quantitative Spectroscopy and Radiative Transfer\\
	extbf{Volume}: 330 (2025), 109222\\[4pt]

	extbf{Penulis}: Favre, Bultel, Morel, Lesage, Gosse
\end{block}

\begin{exampleblock}{Masalah yang Diselesaikan}
\small
Rekonstruksi spektrum LIBS berbasis model fisika untuk sampel kompleks tanpa memerlukan parameter kalibrasi empiris.
\end{exampleblock}

\end{column}

\begin{column}{0.40\textwidth}

\begin{block}{Kontribusi Utama}
\small
\begin{itemize}
\item Model radiative transfer LTE adaptif
\item Simulasi spektrum plasma berbasis Saha--Boltzmann
\item Integrasi basis data atomik otomatis
\item Validasi pada material Eurofer97
\end{itemize}
\end{block}

\begin{alertblock}{Inti Pesan}
\small
Spektrum LIBS dapat direkonstruksi secara konsisten menggunakan model fisika plasma tanpa parameter bebas.
\end{alertblock}

\end{column}

\end{columns}

\end{frame}


% ------------------------------------------------------------
\begin{frame}{Paper Favre 2025: Machine Learning untuk CF-LIBS}

\begin{columns}[T]

\begin{column}{0.56\textwidth}

\begin{block}{Identitas Artikel}
\small
	extbf{Judul}: \textit{Towards real-time calibration-free LIBS supported by machine learning}\\[4pt]

	extbf{Jurnal}: Spectrochimica Acta Part B\\
	extbf{Tahun}: 2025\\[4pt]

	extbf{Penulis}: Favre et al.
\end{block}

\begin{exampleblock}{Masalah yang Diselesaikan}
\small
CF-LIBS klasik memerlukan proses analitik yang kompleks untuk menentukan parameter plasma dan konsentrasi unsur.
\end{exampleblock}

\end{column}

\begin{column}{0.40\textwidth}

\begin{block}{Kontribusi Utama}
\small
\begin{itemize}
\item Dataset simulasi spektrum berbasis fisika plasma
\item Model CNN untuk prediksi komposisi unsur
\item Estimasi parameter plasma secara langsung
\item Implementasi analisis LIBS real-time
\end{itemize}
\end{block}

\begin{alertblock}{Inti Pesan}
\small
Deep learning dapat menggantikan sebagian proses analitik CF-LIBS untuk estimasi parameter plasma dan komposisi unsur.
\end{alertblock}

\end{column}

\end{columns}

\end{frame}


% ------------------------------------------------------------
\begin{frame}{Paper Liu 2026: Transformer-CNN untuk LIBS}

\begin{columns}[T]

\begin{column}{0.56\textwidth}

\begin{block}{Identitas Artikel}
\small
	extbf{Judul}: \textit{Remote LIBS based on transformer-CNN method for quantitative analysis of trace elements in steel}\\[4pt]

	extbf{Jurnal}: IOP Conference Series\\
	extbf{Tahun}: 2026\\[4pt]

	extbf{Penulis}: Liu et al.
\end{block}

\begin{exampleblock}{Masalah yang Diselesaikan}
\small
Analisis unsur jejak pada LIBS jarak jauh sulit karena sinyal lemah dan interferensi garis spektrum.
\end{exampleblock}

\end{column}

\begin{column}{0.40\textwidth}

\begin{block}{Kontribusi Utama}
\small
\begin{itemize}
\item Model hibrida Transformer--CNN
\item CNN mengekstraksi fitur lokal spektrum
\item Transformer mempelajari dependensi global
\item Peningkatan akurasi analisis unsur jejak
\end{itemize}
\end{block}

\begin{alertblock}{Inti Pesan}
\small
Kombinasi CNN dan Transformer meningkatkan kemampuan model dalam memahami struktur kompleks spektrum LIBS.
\end{alertblock}

\end{column}

\end{columns}

\end{frame}


% ------------------------------------------------------------
\begin{frame}{Paper Walidain 2026: Informer untuk Analisis LIBS}

\begin{columns}[T]

\begin{column}{0.56\textwidth}

\begin{block}{Identitas Artikel}
\small
	extbf{Judul}: \textit{Informer-Based LIBS for Qualitative Multi-Element Analysis of an Aceh Traditional Women’s Medicine}\\[4pt]

	extbf{Konferensi}: IOP Conference Series\\
	extbf{Tahun}: 2026\\[4pt]

	extbf{Penulis}: Walidain et al.
\end{block}

\begin{exampleblock}{Masalah yang Diselesaikan}
\small
Analisis multi-elemen pada spektrum LIBS memerlukan metode yang mampu memahami hubungan kompleks antar fitur spektrum.
\end{exampleblock}

\end{column}

\begin{column}{0.40\textwidth}

\begin{block}{Kontribusi Utama}
\small
\begin{itemize}
\item Implementasi Informer Transformer
\item Analisis kualitatif multi-elemen
\item Pemanfaatan self-attention pada spektrum
\item Peningkatan interpretasi pola spektrum
\end{itemize}
\end{block}

\begin{alertblock}{Inti Pesan}
\small
Transformer berbasis attention efektif untuk memahami hubungan kompleks antar garis spektrum LIBS.
\end{alertblock}

\end{column}

\end{columns}

\end{frame}

\end{document}